\chapter{Quick-Finder Guide: Matching Symptoms and Body Areas to Scans}
\label{app:quick-finder}

\noindent
This quick-finder helps you locate the most relevant chapters based on the body area or symptoms you are experiencing. Because doctors choose tests using your full medical history, urgency, age, pregnancy status, kidney health, and implants, use this guide to understand what your care team may consider---never as a substitute for medical advice.

\section{Brain, Nerves, and Spine}
{\renewcommand{\arraystretch}{1.25}\noindent
\begin{tabularx}{\textwidth}{@{}p{0.31\textwidth} p{0.33\textwidth} X@{}}
\toprule
\textbf{Common Symptom or Concern} & \textbf{Typical First-Line Scan} & \textbf{Advanced / Follow-Up Options} \\
\midrule
Sudden stroke signs (face droop, arm weakness, slurred speech) & Head CT (Ch.~\ref{ch:infographic-ctscan}), CTA (Ch.~\ref{ch:infographic-cta}), CT Perfusion (Ch.~\ref{ch:infographic-ctp}) & MRI (Ch.~\ref{ch:infographic-mri}), DWI (Ch.~\ref{ch:infographic-dwi}), PWI (Ch.~\ref{ch:infographic-pwi}), DSA (Ch.~\ref{ch:infographic-dsa}) \\
Severe thunderclap headache or suspected aneurysm & Head CT (Ch.~\ref{ch:infographic-ctscan}), CTA (Ch.~\ref{ch:infographic-cta}) & MRA (Ch.~\ref{ch:infographic-mra}), MRV (Ch.~\ref{ch:infographic-mrv}), DSA (Ch.~\ref{ch:infographic-dsa}) \\
Head injury or lingering post-concussion symptoms & Head CT (Ch.~\ref{ch:infographic-ctscan}) & MRI (Ch.~\ref{ch:infographic-mri}), SWI (Ch.~\ref{ch:infographic-swi}), DTI (Ch.~\ref{ch:infographic-sti}) \\
Seizures or medication-resistant epilepsy & MRI (Ch.~\ref{ch:infographic-mri}) & PET (Ch.~\ref{ch:infographic-petscan}), SPECT (Ch.~\ref{ch:infographic-spect}), MEG (Ch.~\ref{ch:infographic-meg}), fMRI (Ch.~\ref{ch:infographic-functional}) \\
Progressive memory loss or tremors & MRI (Ch.~\ref{ch:infographic-mri}) & PET (Ch.~\ref{ch:infographic-petscan}), SPECT / DaTscan (Ch.~\ref{ch:infographic-spect}) \\
Brain tumor mapping before surgery & MRI (Ch.~\ref{ch:infographic-mri}) & fMRI (Ch.~\ref{ch:infographic-functional}), Fiber Tractography (Ch.~\ref{ch:infographic-ft}), MRS (Ch.~\ref{ch:infographic-mrs}), MEG (Ch.~\ref{ch:infographic-meg}) \\
Back or neck pain shooting down an arm or leg (sciatica) & MRI (Ch.~\ref{ch:infographic-mri}) or X-Ray (Ch.~\ref{ch:infographic-xray}) & Myelography (Ch.~\ref{ch:infographic-myelography}) if MRI is not possible; SPECT/CT (Ch.~\ref{ch:infographic-spect-ct}) \\
\bottomrule
\end{tabularx}}

\section{Chest, Lungs, and Heart}
{\renewcommand{\arraystretch}{1.25}\noindent
\begin{tabularx}{\textwidth}{@{}p{0.31\textwidth} p{0.33\textwidth} X@{}}
\toprule
\textbf{Common Symptom or Concern} & \textbf{Typical First-Line Scan} & \textbf{Advanced / Follow-Up Options} \\
\midrule
Persistent cough, fever, or suspected pneumonia & Chest X-Ray (Ch.~\ref{ch:infographic-xray}) & Chest CT (Ch.~\ref{ch:infographic-ctscan}), POCUS (Ch.~\ref{ch:infographic-pocus}), DEI (Ch.~\ref{ch:infographic-dei}) \\
Sudden sharp chest pain or collapsed lung & Chest X-Ray (Ch.~\ref{ch:infographic-xray}), POCUS (Ch.~\ref{ch:infographic-pocus}) & Chest CT (Ch.~\ref{ch:infographic-ctscan}), DEI (Ch.~\ref{ch:infographic-dei}) \\
Suspected lung blood clot (pulmonary embolism) & CTA (Ch.~\ref{ch:infographic-cta}) & Nuclear V/Q Scintigraphy (Ch.~\ref{ch:infographic-nms}) if avoiding CT dye \\
Chest pressure or checking heart arteries & Cardiac CT (Ch.~\ref{ch:infographic-cardiac-ct}) or Nuclear Stress SPECT (Ch.~\ref{ch:infographic-spect}) & PET (Ch.~\ref{ch:infographic-petscan}), DSA / Angiography (Ch.~\ref{ch:infographic-dsa}) \\
Shortness of breath, leg swelling, or heart murmur & Echocardiography (Ch.~\ref{ch:infographic-echocardiography}), Chest X-Ray (Ch.~\ref{ch:infographic-xray}) & Cardiac MRI (Ch.~\ref{ch:infographic-cardiac-mri}), Cardiac CT (Ch.~\ref{ch:infographic-cardiac-ct}), CEUS (Ch.~\ref{ch:infographic-ceus}) \\
\bottomrule
\end{tabularx}}

\section{Abdomen, Liver, and Digestive Tract}
{\renewcommand{\arraystretch}{1.25}\noindent
\begin{tabularx}{\textwidth}{@{}p{0.31\textwidth} p{0.33\textwidth} X@{}}
\toprule
\textbf{Common Symptom or Concern} & \textbf{Typical First-Line Scan} & \textbf{Advanced / Follow-Up Options} \\
\midrule
Upper-right belly pain after meals (gallstones) & Ultrasound (Ch.~\ref{ch:infographic-ultrasound}) & MRCP (Ch.~\ref{ch:infographic-mrcp}), HIDA Scintigraphy (Ch.~\ref{ch:infographic-nms}) \\
Yellowing eyes or skin (jaundice) & Ultrasound (Ch.~\ref{ch:infographic-ultrasound}) & MRCP (Ch.~\ref{ch:infographic-mrcp}), EUS (Ch.~\ref{ch:infographic-eus}), CT (Ch.~\ref{ch:infographic-ctscan}) \\
Fatty liver or checking liver scarring & Ultrasound Elastography (Ch.~\ref{ch:infographic-ultrasound-elastography}) & MR Elastography (Ch.~\ref{ch:infographic-mre}), MRI (Ch.~\ref{ch:infographic-mri}), CEUS (Ch.~\ref{ch:infographic-ceus}) \\
Sudden severe belly or flank pain (appendicitis, kidney stone) & CT (Ch.~\ref{ch:infographic-ctscan}) or Ultrasound (Ch.~\ref{ch:infographic-ultrasound}) & Dual-Energy CT (Ch.~\ref{ch:infographic-dual-energy-ct}) for stone type; MRI (Ch.~\ref{ch:infographic-mri}) in pregnancy; POCUS (Ch.~\ref{ch:infographic-pocus}) \\
Pregnancy checkups, fetal growth, or prenatal anatomy & Prenatal Ultrasound (Ch.~\ref{ch:infographic-obstetric-ultrasound}) & Doppler Ultrasound (Ch.~\ref{ch:infographic-doppler-ultrasound}), Fetal MRI (Ch.~\ref{ch:infographic-mri}) \\
Trouble swallowing, reflux, or upper-belly pain & Endoscopy (Ch.~\ref{ch:infographic-endoscopy}) & Fluoroscopy / Barium Swallow (Ch.~\ref{ch:infographic-fluroscopy}), EUS (Ch.~\ref{ch:infographic-eus}) \\
Unexplained anemia, hidden gut bleeding, or bowel screening & Endoscopy / Colonoscopy (Ch.~\ref{ch:infographic-endoscopy}) & Capsule Endoscopy (Ch.~\ref{ch:infographic-capsule-endoscopy}), CTA (Ch.~\ref{ch:infographic-cta}), IR (Ch.~\ref{ch:infographic-interventional-radiology}) \\
\bottomrule
\end{tabularx}}

\section{Bones, Joints, Teeth, and Jaw}
{\renewcommand{\arraystretch}{1.25}\noindent
\begin{tabularx}{\textwidth}{@{}p{0.31\textwidth} p{0.33\textwidth} X@{}}
\toprule
\textbf{Common Symptom or Concern} & \textbf{Typical First-Line Scan} & \textbf{Advanced / Follow-Up Options} \\
\midrule
Suspected broken bone, dislocation, or arthritis & X-Ray (Ch.~\ref{ch:infographic-xray}) & CT (Ch.~\ref{ch:infographic-ctscan}) for complex fractures; MRI (Ch.~\ref{ch:infographic-mri}) for hidden cracks \\
Torn knee/shoulder ligament, cartilage, or tendon & MRI (Ch.~\ref{ch:infographic-mri}) or Ultrasound (Ch.~\ref{ch:infographic-ultrasound}) & Fluoroscopy-guided arthrogram (Ch.~\ref{ch:infographic-fluroscopy}) \\
Osteoporosis screening, height loss, or fragility fracture & DEXA (Ch.~\ref{ch:infographic-dexa}) & Spine X-Ray (Ch.~\ref{ch:infographic-xray}) or MRI (Ch.~\ref{ch:infographic-mri}) \\
Toothache, cavities, wisdom teeth, or jaw injury & Dental Radiography (Ch.~\ref{ch:infographic-dental-radiography}), OPG (Ch.~\ref{ch:infographic-opg}) & Dental CBCT (Ch.~\ref{ch:infographic-cbct}) for 3D roots and implants \\
Persistent pain around surgical screws or artificial joints & X-Ray (Ch.~\ref{ch:infographic-xray}) & SPECT/CT (Ch.~\ref{ch:infographic-spect-ct}), Bone Scintigraphy (Ch.~\ref{ch:infographic-nms}) \\
\bottomrule
\end{tabularx}}

\section{Breast Health and Blood Vessels}
{\renewcommand{\arraystretch}{1.25}\noindent
\begin{tabularx}{\textwidth}{@{}p{0.31\textwidth} p{0.33\textwidth} X@{}}
\toprule
\textbf{Common Symptom or Concern} & \textbf{Typical First-Line Scan} & \textbf{Advanced / Follow-Up Options} \\
\midrule
Routine breast screening or dense breasts & Mammography (Ch.~\ref{ch:infographic-mamography}), 3D DBT (Ch.~\ref{ch:infographic-dbt}) & Breast Ultrasound (Ch.~\ref{ch:infographic-breast-ultrasound}), Breast MRI (Ch.~\ref{ch:infographic-breast-mri}) \\
New breast lump, focal pain, or nipple discharge & Mammography (Ch.~\ref{ch:infographic-mamography}) / DBT (Ch.~\ref{ch:infographic-dbt}) + Breast Ultrasound (Ch.~\ref{ch:infographic-breast-ultrasound}) & Breast MRI (Ch.~\ref{ch:infographic-breast-mri}), Ultrasound Elastography (Ch.~\ref{ch:infographic-ultrasound-elastography}) \\
One-sided calf swelling or pain (suspected blood clot / DVT) & Doppler Ultrasound (Ch.~\ref{ch:infographic-doppler-ultrasound}) & MRV (Ch.~\ref{ch:infographic-mrv}), CTA (Ch.~\ref{ch:infographic-cta}) \\
Leg cramping when walking or narrowed neck arteries & Doppler Ultrasound (Ch.~\ref{ch:infographic-doppler-ultrasound}) & CTA (Ch.~\ref{ch:infographic-cta}), MRA (Ch.~\ref{ch:infographic-mra}), DSA (Ch.~\ref{ch:infographic-dsa}) \\
Active internal bleeding or blocked vessel needing repair & CTA (Ch.~\ref{ch:infographic-cta}) & Interventional Radiology (Ch.~\ref{ch:infographic-interventional-radiology}), DSA (Ch.~\ref{ch:infographic-dsa}) \\
\bottomrule
\end{tabularx}}

\section{Eyes and Vision}
{\renewcommand{\arraystretch}{1.25}\noindent
\begin{tabularx}{\textwidth}{@{}p{0.31\textwidth} p{0.33\textwidth} X@{}}
\toprule
\textbf{Common Symptom or Concern} & \textbf{Typical First-Line Scan} & \textbf{Advanced / Follow-Up Options} \\
\midrule
Red, gritty, or painful eye; cataracts; routine exam & Slit-Lamp Exam (Ch.~\ref{ch:infographic-slit-lamp}) & Confocal Microscopy (Ch.~\ref{ch:infographic-confocal}) for stubborn infections \\
Diabetic eye check or optic nerve monitoring & Fundus Photography (Ch.~\ref{ch:infographic-fundus}), OCT (Ch.~\ref{ch:infographic-oct}) & OCTA (Ch.~\ref{ch:infographic-octa}), FFA (Ch.~\ref{ch:infographic-ffa}) \\
Wavy straight lines or blurry central vision & OCT (Ch.~\ref{ch:infographic-oct}) & OCTA (Ch.~\ref{ch:infographic-octa}), FFA (Ch.~\ref{ch:infographic-ffa}), ICGA (Ch.~\ref{ch:infographic-icga}) \\
Sudden flashes/floaters behind a dense cataract or blood & B-Scan Ultrasound (Ch.~\ref{ch:infographic-bscan}) & OCT (Ch.~\ref{ch:infographic-oct}) once view clears \\
High eye pressure, eye pain with halos, or behind-iris mass & Slit-Lamp (Ch.~\ref{ch:infographic-slit-lamp}), OCT (Ch.~\ref{ch:infographic-oct}) & Ultrasound Biomicroscopy (Ch.~\ref{ch:infographic-ubm}) \\
Worsening astigmatism, LASIK planning, or morning corneal haze & Corneal Topography (Ch.~\ref{ch:infographic-corneal-topography}), Corneal Tomography (Ch.~\ref{ch:infographic-corneal-tomography}) & Specular Microscopy (Ch.~\ref{ch:infographic-specular-microscopy}) \\
\bottomrule
\end{tabularx}}

\section{Cancer Staging, Whole-Body Checks, and Research-Stage Scans}
{\renewcommand{\arraystretch}{1.25}\noindent
\begin{tabularx}{\textwidth}{@{}p{0.31\textwidth} p{0.33\textwidth} X@{}}
\toprule
\textbf{Common Symptom or Concern} & \textbf{Typical First-Line Scan} & \textbf{Advanced / Follow-Up Options} \\
\midrule
Whole-body cancer staging or checking treatment response & CT (Ch.~\ref{ch:infographic-ctscan}), MRI (Ch.~\ref{ch:infographic-mri}) & PET/CT (Ch.~\ref{ch:infographic-pet-ct}), PET/MRI (Ch.~\ref{ch:infographic-pet-mri}), PET (Ch.~\ref{ch:infographic-petscan}), SPECT/CT (Ch.~\ref{ch:infographic-spect-ct}) \\
Changing, irregular, or suspicious skin moles & Dermoscopy (Ch.~\ref{ch:infographic-dermoscopy}) & Confocal Microscopy (Ch.~\ref{ch:infographic-confocal}) \\
Overactive thyroid or high calcium (parathyroid) & Ultrasound (Ch.~\ref{ch:infographic-ultrasound}), Scintigraphy (Ch.~\ref{ch:infographic-nms}) & SPECT/CT (Ch.~\ref{ch:infographic-spect-ct}) \\
Chronic diabetic foot wounds or surgical tissue blood flow & Doppler Ultrasound (Ch.~\ref{ch:infographic-doppler-ultrasound}) & Hyperspectral Imaging (Ch.~\ref{ch:infographic-hsi}) at specialized centers \\
Investigational / Research-Stage Technologies & \emph{Not for routine clinical care} & MPI (Ch.~\ref{ch:infographic-mpi}), Raman (Ch.~\ref{ch:infographic-raman}), OAT (Ch.~\ref{ch:infographic-oat}), PAI (Ch.~\ref{ch:infographic-pai}), EIT (Ch.~\ref{ch:infographic-eit}), Microwave (Ch.~\ref{ch:infographic-microwave-imaging}), Terahertz (Ch.~\ref{ch:infographic-terahertz-imaging}) \\
\bottomrule
\end{tabularx}}
